\begin{abstract}
Web Application Firewalls (WAFs) are a frontline defense against SQL injection (SQLi), yet systematic diagnosis of their policy coverage gaps remains methodologically underspecified. Prior adversarial work---including reinforcement-learning agents such as SSQLi reporting bypass rates up to 97.39\% \cite{guan2023ssqli}, and discrepancy-based HTTP mutations \cite{akhavani2025waffled}---demonstrates that policy failures are real and measurable, but rarely decomposes which mutation families expose genuine coverage holes versus which are effectively defended. Without such decomposition, security teams cannot prioritize rule improvements or validate that a patch closes a specific gap.

This work adopts a \emph{defensive audit} perspective and reports a toolchain that (i) composes dialect seeds (PostgreSQL- and Oracle-oriented), deployment-specific augmentations, and learned transform rules into a large discrete variant set for systematic WAF coverage testing; (ii) issues black-box HTTP probes against a live AWS WAF-fronted application with per-attempt JSONL logging; and (iii) aggregates outcomes globally and by mutation \emph{family}, enabling defenders to identify high-risk policy gaps versus well-defended operator classes. On the archived last complete live run (\texttt{aws\_probe\_summary.json}: $N{=}5{,}000$ variants tested), the audit identifies $N_{\mathrm{gap}}{=}2{,}083$ attempts (41.7\%) as policy gaps (WAF-pass plus application-visible signals under the script's definition). Family-level analysis reveals critical heterogeneity with direct remediation implications: identity-, case-mix-, and null-byte-oriented spawn cohorts expose near-complete coverage gaps ($\mathrm{BR}_f{\approx}99\%$) in that run, while encoding- and hex-style spawn cohorts are fully defended ($\mathrm{BR}_f{=}0\%$), indicating that defensive policy improvement should concentrate on a specific subset of syntactic strategies rather than treating all mutations uniformly.

As a secondary control, we apply repeated RFC percent-decoding to a fixpoint before scoring strings with a small TF-IDF + logistic-regression surrogate (\texttt{paper-defense-rfc-ml}); test accuracy stays at 42.9\% for both raw and canonicalized inputs on a 20-request pilot (10 attack-like and 10 benign-like HTTP templates), demonstrating that naive single-pass normalization is insufficient as a standalone defense measure---reinforcing the need for layered defenses addressing richer feature interactions. The methodology's value lies in its explicit decomposition, reproducible artifacts, and family-stratified coverage analysis that enables targeted policy hardening and regression testing after rule updates.
\end{abstract}
